\documentclass{article}
\usepackage{amsmath} 
\usepackage{amsfonts}
\usepackage{amssymb}
\usepackage[parfill]{parskip}
\usepackage{esvect}
\usepackage[thinc]{esdiff}
\usepackage{mathtools}
\usepackage{physics}
\usepackage{amsthm}
\usepackage{xfrac}
\usepackage{enumitem}

\DeclareMathOperator{\lub}{lub}
\DeclareMathOperator{\gub}{gub}

\begin{document}

\title{MATH 411: Week 6}
\author{Jacob Lockard}
\date{2 March 2026}

\maketitle

\section*{Problem 8.9}

\begin{quote}
    \itshape
    Suppose that $f$ and $g$ are functions defined on $[a,b]$ and are both
    differentiable at a point $x \in [a,b]$. Show that if $g(x) \neq 0$,
    \begin{align*}
        \Bigl( \frac{f}{g} \Bigr)' (s) = \frac{ f'(x)g(x) - f(x)g'(x) }{ (g(x))^2} \,.
    \end{align*}
\end{quote}

Let $f$ and $g$ be defined on $[a,b]$ and differentiable at some $x \in [a,b]$,
with $g(x) \neq 0$. We'll show that
\begin{align*}
    \Bigl( \frac{f}{g} \Bigr)' (s) = \frac{ f'(x)g(x) - f(x)g'(x) }{ (g(x))^2} \,.
\end{align*}
We first observe that $\bigl(g^{-1}\bigr)'(x)$ is given by
\begin{align*}
    \bigl(g^{-1}\bigr)'(x) &= \lim_{t\to x} \frac{\frac{1}{g(t)} - \frac{1}{g(x)}}{t-x} \\
    &= \lim_{t\to x} \frac{\frac{g(x) - g(t)}{g(t)g(x)}}{t-x} \\
    &= \lim_{t\to x} \Biggl( \frac{1}{g(t)g(x)} \cdot \frac{g(x) - g(t)}{t-x} \Biggr) \\
    &= \lim_{t\to x} \frac{1}{g(t)g(x)} \cdot \lim_{t\to x} \frac{g(x) - g(t)}{t-x} \\
    &= \frac{1}{g(x)^2}(-g'(x)) \\
    &= -\frac{g'(x)}{g(x)^2} \,,
\end{align*}
since Problem 8.6 implies $g$ is continuous at $x$ and hence that
$t \mapsto \frac{1}{g(t)g(x)}$ is continuous at $x$, since $g(x) \neq 0$
Problem 8.8 now ensures
\begin{align*}
    \Bigl( \frac{f}{g} \Bigr)'(x) = ( fg^{-1} )'(x) &= 
    f(x)\bigl(g^{-1}\bigr)'(x) + f'(x)g^{-1}(x) \\
    &= -f(x) \frac{g'(x)}{g(x)^2} + f'(x)\frac{1}{g(x)} \\
    &= -f(x) \frac{g'(x)}{g(x)^2} + f'(x)\frac{g(x)}{g(x)^2} \\
    &= \frac{f'(x)g(x) - f(x)g'(x)}{g(x)^2} \,.
\end{align*}

\qedsymbol

\section*{Problem 8.10}

\begin{quote}
    \itshape
    Suppose that $f$ is continous on $[a,b]$, $f$ is differentiable at some
    point $x \in [a,b]$, $g$ is defined on an interval which contains
    the range of $f$, and $g$ is differentiable at $f(x)$. Let
    \begin{align*}
        h(t) = g(f(t)) \quad \text{for } t \in [a,b] \,.
    \end{align*}
    Prove that $h$ is differentiable at $x$ and
    \begin{align*}
        h'(x) = g'(f(x))f'(x) \,.
    \end{align*}
\end{quote}

Let $f$ be continous on $[a,b]$ and differentiable at some $x \in [a,b]$.
Let $g$ be defined on some interval $[\alpha, \beta]$ that 

By definition,
\begin{align*}
    g'(f(x)) = \lim_{t \to f(x)} \frac{g(t) - g(f(x))}{t-f(x)} \,.
\end{align*}
This means that there's

\section*{Problem 8.13}

Let $f$ and $g$ be continuous real-valued functions on $[a,b]$
which are differentiable on $(a,b)$. Define, for all $t  \in [a,b]$,
\begin{align*}
    h(t) = [f(b) - f(a)]g(t) - [g(b) - g(a)]f(t) \,.
\end{align*}

We first show that $h$ is continuous on $[a,b]$ and differentiable on $(a,b)$.
Problem 6.6 ensures $h$ is continuous on $[a,b]$.
$t \mapsto f(b) - f(a)$, $t \mapsto g(b) - g(a)$, $g$, and $f$ are all differentiable
on $(a,b)$, so Problem 8.8 and 8.7 ensures $h$ is also differentiable on $(a,b)$.

Now we show that $h(a) = h(b)$:
\begin{align*}
    h(a) &= [f(b) - f(a)]g(a) - [g(b) - g(a)]f(a) \\
    &= f(b)g(a) - f(a)g(a) - [g(b)f(a) - g(a)f(a)] \\
    &= f(b)g(a) - f(a)g(a) - g(b)f(a) + g(a)f(a) \\
    &= f(b)g(a) - g(b)f(a) \\
    &= g(b)f(b) - g(b)f(a) - g(b)f(b) + f(b)g(a)  \\
    &= [f(b) - f(a)]g(b) - [g(b) + g(a)]f(b)  \\
    &= h(b) \,.
\end{align*}

Finally, suppose $h'(x) = 0$ for some $x \in (a,b)$. We'll show that
\begin{align*}
    [f(b) - f(a)]g'(x) = [g(b) - g(a)]f'(x) \,.
\end{align*}
Since $t \mapsto f(b) - f(a)$ and $t \mapsto g(b) -g(a)$ are constant functions,
Problems 8.7 and 8.8 ensure
\begin{align*}
    h'(t) = [f(b)-f(a)]g'(t) - [g(b)-g(a)]f'(t) \,.
\end{align*}
So we have:
\begin{align*}
    0 = h'(x) &= [f(b) - f(a)]g'(x) - [g(b) - g(a)]f'(x) \\
    [f(b) - f(a)]g'(x) &= [g(b) - g(a)]f'(x) \,.
\end{align*}

\qedsymbol

\section*{Problem 8.14}

\begin{quote}
    \itshape
    Using the conclusions from the previous problem, show that there exists
    some $x \in (a,b)$ such that $h'(x) = 0$.
\end{quote}

Problem 6.13 ensures that $h$ attains a maximum value $M$ and a minimum value $m$
at some $x_{\text{max}}$ and $x_{\text{min}}$, respectively, in $[a,b]$.

If either $x_{\text{min}} \in (a,b)$ or $x_{\text{max}} \in (a,b)$, then
Theorem 8.1 ensures either $f'(x_{\text{min}}) = 0$ or $f'(x_{\text{max}}) = 0$.
On the other hand, if $x_{\text{min}} = a$ and $x_{\text{max}} = b$, then we have
\begin{align*}
    M = f(x_{\text{max}}) = f(a) = f(b) = f(x_{\text{min}}) = m \,,
\end{align*}
and if $x_{\text{max}} = b$ and $x_{\text{min}} = a$, then we have
\begin{align*}
    m = f(x_{\text{min}}) = f(a) = f(b) = f(x_{\text{max}}) = M \,.
\end{align*}
But if $m = M$, then for all $x \in [a,b]$ we have
$h(x) \geq m = M$ and $h(x) \leq M$; so $h(x) = M$ and $h'(x) = 0$ for any $x \in [a,b]$.
In any case, we've shown that an $x$ with the desired properties exist.

\qedsymbol

\end{document}